% ======================================================================
% SECTION 11: SITE INTERFACE LAYER AND INVESTIGATOR COORDINATION
% File: sections/site_interface.tex
% ======================================================================

The site interface layer manages all interactions between the autonomous sponsor and the
physical trial sites where patients are enrolled and treated. Traditional sponsors interact
with sites through clinical research associates, medical monitors, and project managers;
the autonomous sponsor replaces these interactions with the site gateway agent, which
provides standardized interfaces for site qualification, investigator coordination, systems
integration, and patient-facing services. The site interface must accommodate both the
conventional clinical trial activities (consent, data collection, monitoring visits) and
the Physical AI-specific activities (robotic system authorization, digital twin
synchronization, and telemetry data exchange).

\subsection{Site Gateway}

The site gateway (site\_gateway) manages all sponsor-to-site interactions through a unified
interface. Site qualification and readiness assessment uses the Physical AI Standard Level
(PSL) framework \cite{national-platform-paper} to evaluate sites across three dimensions:
Legislative compliance (conformance with legislation authorization, patient rights, and
data transparency requirements), Regulatory compliance (conformance with building codes,
premises codes, parking and transportation requirements, site operations, and emergency
preparedness standards), and Operational readiness (staff competency, system availability,
and infrastructure capability).

The 11-document site establishment package \cite{site-docs} provides the foundation for
site readiness assessment. Each document category maps to specific PSL scoring criteria:
legislation authorization, patient rights legislation, data transparency legislation,
city regulations, state regulations, national regulations (FDA), building code, premises
code, parking and transportation, site operations, and emergency preparedness. Sites must
achieve passing scores across all PSL dimensions before activation for trial participation.

Training module delivery automates the investigator and site staff training required before
trial participation. The gateway generates role-specific training curricula: investigators
receive protocol training, GCP training, and Physical AI system orientation; research
coordinators receive protocol procedures training, EDC training, and safety reporting
training; pharmacy staff receive IMP handling and robotic pharmacy interface training;
nursing staff receive patient care procedures and decentralized visit support training.
Competency assessments verify training effectiveness, with automated retraining triggers
when competency scores fall below thresholds.

Scheduling and resource allocation optimizes the use of limited trial site resources:
robotic system time, treatment room availability, imaging equipment, pharmacy preparation
capacity, and clinical staff availability. The gateway generates optimized schedules that
minimize patient wait times while maximizing resource utilization. For the 24-hour
simulation, this scheduling approach achieved an average patient wait time of 12 minutes
from arrival to procedure start across 168 patients \cite{main-repo}.

\subsection{Investigator Interface}

The investigator interface provides the principal investigator and sub-investigators with
visibility into and control over the trial activities at their site. Protocol training and
certification automation delivers interactive protocol training, assesses comprehension
through structured assessments, and maintains certification records. Annual recertification
is triggered automatically with updated content reflecting protocol amendments and new
safety information.

Delegation of authority log management maintains the electronic record of which trial
activities each investigator and qualified staff member is authorized to perform.
Delegation updates are processed through the site gateway with electronic signatures from
the principal investigator. The system prevents task assignment to staff without current
delegation authority and training certification.

Medical monitoring communication channels provide structured pathways for investigators to
consult with the clinical accountability agent on medical questions arising during trial
conduct. Urgent clinical questions (patient safety, eligibility exceptions, dose
modifications) receive immediate routing with defined response timelines. Non-urgent
questions enter a standard queue with periodic batch review.

The investigator oversight dashboard provides real-time visibility into site-level trial
metrics: enrollment status, data quality indicators, safety event summary, query status,
and monitoring visit findings. The dashboard enables investigators to exercise their
ICH E6(R3) responsibility for the conduct of the trial at their site while the autonomous
sponsor manages operational coordination.

\subsection{Site Systems Integration}

Site systems integration connects the autonomous sponsor with the clinical information
systems operating at each trial site. Table~\ref{tab:site-systems} presents the
integration matrix.

\begin{longtable}{L{1.8cm} L{1.5cm} L{2.0cm} L{1.5cm} L{2.5cm} L{2.2cm}}
\caption{Site Systems Integration Matrix}
\label{tab:site-systems} \\
\toprule
\textbf{System} & \textbf{Protocol} & \textbf{Data Type} & \textbf{Direction} & \textbf{Agent Handler} & \textbf{Security} \\
\midrule
\endfirsthead
\toprule
\textbf{System} & \textbf{Protocol} & \textbf{Data Type} & \textbf{Direction} & \textbf{Agent Handler} & \textbf{Security} \\
\midrule
\endhead
\midrule
\multicolumn{6}{r}{\textit{Continued on next page}} \\
\endfoot
\bottomrule
\endlastfoot
EHR & FHIR R4 & Demographics, history, labs, vitals, medications & Bidirectional & data\_biostats\_agent, site\_gateway & TLS 1.3, OAuth 2.0, SMART on FHIR \\
PACS & DICOM & Imaging studies (CT, MRI, PET, X-ray) & Inbound & data\_biostats\_agent & TLS 1.3, DICOM TLS \\
LIS & HL7 v2 & Laboratory orders and results & Bidirectional & data\_biostats\_agent & TLS 1.3, HL7 security \\
Pharmacy & Custom API & IMP inventory, orders, dispensing records & Bidirectional & supply\_agent & TLS 1.3, mutual TLS \\
MCP Server & MCP protocol & Structured tool invocations and responses & Bidirectional & All agents via MCP client & MCP authentication \\
Robot Systems & ROS/API & Telemetry, commands, status & Bidirectional & robot\_execution\_gateway & TLS 1.3, certificate pinning \\
\end{longtable}

Table~\ref{tab:site-readiness} presents the site readiness assessment criteria.

\begin{longtable}{L{2.0cm} L{3.0cm} L{2.2cm} L{1.5cm} L{3.0cm}}
\caption{Site Readiness Assessment Criteria (PSL-Based)}
\label{tab:site-readiness} \\
\toprule
\textbf{Category} & \textbf{Criterion} & \textbf{Scoring Method} & \textbf{Minimum} & \textbf{Evidence Required} \\
\midrule
\endfirsthead
\toprule
\textbf{Category} & \textbf{Criterion} & \textbf{Scoring Method} & \textbf{Minimum} & \textbf{Evidence Required} \\
\midrule
\endhead
\midrule
\multicolumn{5}{r}{\textit{Continued on next page}} \\
\endfoot
\bottomrule
\endlastfoot
Legislative & Trial authorization and patient rights compliance & Pass/fail per PSL legislative dimension & Pass all & Signed legislation compliance attestation \\
Regulatory & Building, premises, and operational code compliance & Pass/fail per PSL regulatory dimension & Pass all & Inspection certificates, compliance documentation \\
Infrastructure & Robotic system installation, calibration, and qualification & USL scoring per robot category & $\geq$ 75.0 & Calibration records, qualification protocols \\
IT Systems & EHR, PACS, LIS connectivity and data exchange validation & Integration testing pass rate & $\geq$ 95\% & Integration test reports, data flow verification \\
Staff & Investigator and staff training and competency & Assessment scores by role & $\geq$ 80\% & Training certificates, competency assessments \\
Safety & Emergency preparedness and safety system validation & Emergency drill completion and safety gate testing & Pass all & Drill reports, safety system validation records \\
\end{longtable}

\subsection{Patient-Facing Interface}

The patient-facing interface manages the trial participant experience from initial contact
through study completion. eConsent platform management delivers the electronic informed
consent process in compliance with FDA eConsent guidance \cite{fda-econsent} and the adapted
21 CFR Part 50 requirements for Physical AI procedures \cite{cfr50-adapt}. The consent
interface presents trial information in an accessible format, provides interactive elements
for patient education about robotic procedures, captures comprehension verification, and
records electronic consent signatures.

Patient scheduling automation coordinates trial visits with the site scheduling system and
the robot execution gateway to ensure that required resources (clinical staff, treatment
rooms, robotic systems, imaging equipment) are available for each patient visit. The
scheduling system manages protocol-defined visit windows, optimizes appointment clustering
to reduce patient travel burden, and accommodates patient preferences where protocol
flexibility permits.

Electronic patient-reported outcome collection deploys validated instruments through
mobile and web platforms. The agent configures instrument deployment schedules aligned
with the protocol assessment windows, monitors completion rates, generates reminders for
overdue assessments, and flags clinically significant responses for clinical review.
Patient-reported safety data integrates with the safety agent (\S\ref{sec:safety}) for
continuous safety monitoring between site visits.
